%% kapa: exam=1 round=37 year=2026 period=1 form=A subject=civil_law subject_ko=민법 seq=13 number=13
%% source: Q-Net 감정평가사 기출문제 (cst00309) artlSeq=5258924 file=raw/1008/5258924/2026년 제37회 감정평가사 1차 1교시.pdf
%% conversion: kordoc 4.24.51 -> tools/gen1.py (auto); review=proofread-sonnet; flags=none
\begin{question}{13}
사기에 의한 의사표시에 관한 설명으로 옳지 않은 것은? (다툼이 있으면 판례에 따름)
\begin{choices}
  \item 기망행위로 인한 동기의 착오의 경우, 표의자는 자신의 의사표시를 사기에 의한 의사표시로서 취소할 수 있다.
  \item 사기에 의한 의사표시의 취소의 효과로 생기는 부당이득반환청구권과 불법행위로 인한 손해배상청구권이 경합하여 병존하는 경우, 채권자는 양 청구권을 중첩적으로 행사할 수 있다.
  \item 매수인이 목적물의 시가를 묵비하여 매도인에게 고지하지 않았더라도, 원칙적으로 위법한 기망행위가 되지 않는다.
  \item 매수인이 매도인을 기망하여 부동산을 매수한 후 제3자에게 저당권을 설정해 준 경우, 특별한 사정이 없는 한 그 제3자는 매수인의 기망사실에 대하여 선의인 것으로 추정된다.
  \item 대리인의 기망행위로 계약을 체결한 상대방은 본인이 대리인의 기망행위에 대해 선의ㆍ무과실이더라도 계약을 취소할 수 있다.
\end{choices}
\end{question}
\answer{2}
